(1) 一维均匀材料的体积热容为 \(c_v>0\)、导热率为 \(\kappa>0\)，热流 \(q(x,t)\) 满足 \(\tau q_t+q=-\kappa T_x\)，热流弛豫时间 \(\tau>0\)；局部能量守恒为 \(c_vT_t+q_x=0\)。消去热流得到温度的双曲型方程。
(2) 对无限均匀介质采用平面波 \(T\propto e^{i(kx-\omega t)}\)，求色散关系及高波数热波速度。
(3) 若边界条件保证 \(Tq\) 的边界功率流为零，构造涉及温度扰动与热流的正定二次 Lyapunov 泛函，并证明其时间导数非正。
(4) 检验 \(\tau\to0^+\) 时回到傅里叶扩散模型以及长波模的衰减特征，说明有限热波速度并不等于载流子或光速。
